% 地球（综述）
% license CCBYSA3
% type Wiki

本文根据 CC-BY-SA 协议转载翻译自维基百科\href{https://en.wikipedia.org/wiki/Earth}{相关文章}。

\begin{figure}[ht]
\centering
\includegraphics[width=6cm]{./figures/5280dfaa768e0550.png}
\caption{“蓝色弹珠”，阿波罗 17 号，1972 年 12 月} \label{fig_Earth_1}
\end{figure}
地球是距太阳的第三颗行星，也是唯一已知能够孕育生命的天体。这得益于地球是一颗海洋世界——在太阳系中，唯一能够维持液态地表水的星球。地球上几乎所有的水都存在于全球性海洋中，覆盖了地球地壳的 70.8\%。剩余的 29.2\% 地壳为陆地，其中大部分集中在地球陆半球的大陆地块上。地球的大多数陆地区域至少具有一定湿度并被植被覆盖，而地球极地沙漠中的大型冰盖所储存的水量超过了地球地下水、湖泊、河流和大气水的总和。地球的地壳由缓慢移动的构造板块组成，彼此相互作用形成山脉、火山与地震。地球拥有液态外核，可产生磁层，能够偏转大部分具有破坏性的太阳风和宇宙射线。

地球拥有一个动态的大气层，这维持了地球的表面环境，并在陨石体及紫外线入射时提供保护。大气层主要由氮气和氧气组成。水汽在大气层中广泛存在，形成覆盖地球大部分区域的云层。水汽是一种温室气体，并与大气中的其他温室气体（特别是二氧化碳 $CO_2$）一道，通过捕获来自太阳光的能量，创造出液态表面水和水汽能够持续存在的条件。这一过程维持了目前 14.76 °C（58.57 °F）的平均地表温度，在正常大气压下，水在该温度呈液态。不同地理区域所捕获能量的差异（如赤道地区获得的阳光多于极地地区）驱动了大气和海洋环流，形成具有不同气候区的全球气候系统，并产生降水等多种天气现象，使碳和氮等组成成分得以循环。

地球呈近似椭球形，周长约为 40,000 公里（24,900 英里）。它是太阳系中密度最高的行星。在四颗类地行星中，地球最大且质量最高。地球距太阳约八光分（1 天文单位），沿轨道绕太阳运行一周约需一年（约 365.25 天）。地球自转一周所需时间略少于一天（约 23 小时 56 分钟）。地球自转轴相对于其绕太阳轨道平面的垂线有倾斜，从而产生四季更替。地球拥有一颗永久天然卫星——月球。月球以 384,400 公里（238,855 英里，约 1.28 光秒）的距离绕地球运行，其直径约为地球的四分之一。月球的引力有助于稳定地球自转轴，产生潮汐，并逐渐减缓地球自转。同样，地球的引力已使月球自转潮汐锁定，月球永远以同一面朝向地球。

地球与太阳系中大多数其他天体一样，约在 45 亿年前由早期太阳系中的气体和尘埃形成。在地球历史的最初十亿年中，海洋形成，随后生命在海洋中出现。生命在全球扩散，并持续改变地球的大气和地表，导致约 20 亿年前发生“大氧化事件”。人类约 30 万年前在非洲出现，随后扩散到地球的每一块大陆。人类依赖地球的生物圈和自然资源而生存，但对地球环境的影响却日益增加。如今，人类对地球气候与生物圈的影响已不可持续，威胁着人类自身及许多其他生命形式的生计，并导致物种广泛灭绝。
